\documentclass{article}
\usepackage[margin=0.8in]{geometry}
\usepackage{amsmath,amssymb,amsthm,mathtools}
\usepackage{mathrsfs,bm,bbm}
\usepackage{graphicx,booktabs,array,multirow,tabularx,longtable}
\usepackage{enumitem}
\usepackage{xcolor}
\usepackage{algorithm,algpseudocode}
\usepackage{subcaption,tikz}
\usepackage{siunitx,cancel,accents}
\usepackage{microtype}
\usepackage[numbers,sort&compress]{natbib}
\usepackage[colorlinks=true,linkcolor=blue,citecolor=blue,urlcolor=blue]{hyperref}
\usepackage{cleveref}
\setlength{\emergencystretch}{2em}
\newtheorem{assumption}{Assumption}
\newtheorem{definition}{Definition}
\newtheorem{theorem}{Theorem}
\newtheorem{lemma}{Lemma}
\newtheorem{proposition}{Proposition}
\newtheorem{openendedquestion}{Open-ended Question}
\newtheorem{conjecture}{Conjecture}
\newcommand{\coreref}[1]{\ref{#1}}

\begin{document}

\sloppy

\section*{eid\_\allowbreak{}pbscm\_\allowbreak{}sourceforest\_\allowbreak{}complete\_\allowbreak{}fiber}

\noindent\textit{Specialization:} v1\par

\paragraph{TL;DR.} Event-count branching models encode causal structure beyond conditional independence, yet their strict all-Poisson form admits genuine direction reversals. The proposal characterizes the complete one-law observational fiber. Degree-four exact-support log-PGF coefficients recover the skeleton and every arrow entering a multi-parent vertex; deleting those arrows leaves anchored and free trees, and the fiber consists exactly of independently rerooting the free trees. A positive source-edge reversal realizes every alternative, while a reverse-topological rational decoder yields the unique parameter point on each graph. Degree four determines the entire strict-model law and is minimal by an explicit triangle collision through degree three. A bounded-cell confidence region is outward-inverted to cover the whole population fiber, and only arrows unanimous over a nonempty outer set are reported.

\section{Project justification}

\paragraph{Gap.} Published PB-SCM theory supplies path-based or local PGF orientation criteria but no necessary-and-sufficient description of every DAG and parameter point generating a fixed strict all-Poisson observational law.

\paragraph{Niche.} The local parent-pair signatures leave a structured residual ambiguity that can be classified globally without faithfulness and can be carried into honest set-valued inference.

\paragraph{Fill.} The source-forest theorem gives the exact fiber, transformations, cardinality, certificates, and rational inverse; the sharp degree-four theorem compresses the full PGF; and outward coefficient inversion provides finite-sample coverage of the unrestricted fiber.

\section{Related work}

Qiao et al. (2024) connect arbitrary-order cumulants to directed-path multiplicity, establish sufficient orientation conditions, and give the bivariate reversal that motivates the source-edge transformation. Xiang et al. (2024) derive the recursive PB-SCM PGF and local skeleton, triangle, and collider signatures; their event-count framing and football and shopping real-data experiments provide motivating use cases, while their global observational fiber remains unstated. Ghassami et al. (2020) separate distribution equivalence from conditional-independence equivalence in linear Gaussian directed graphs and provide a methodological template for explicit same-law transformations, but their algebra and equivalence target differ. Gao, Cai, and Hara (2026) show that the all-Poisson PB-SCM is exceptional and obtain full identification after changing the thinning or nonsink innovation model; they do not characterize the exceptional model's fiber. Qiao et al. (2026) address local three-variable monomial equivalence under latent confounding, a different model and scope. Strieder et al. (2021) and Wang, Kolar, and Drton (2026) show how inversion and relation envelopes retain causal-structure uncertainty, but their guarantees concern restricted linear or additive-error models rather than finite-sample coverage of a non-singleton PB-SCM population fiber. Nakamura and Perez-Abreu (1993) provide classical empirical-PGF inference; the present low-degree rational cell map specializes that idea to structural-fiber coverage. Steutel and van Harn (1979) develop discrete self-decomposability and stability, and Al-Osh and Alzaid (1987) introduce a univariate binomial-thinning autoregression; both supply historical count-process context, but neither treats causally sufficient multivariate DAG PB-SCMs nor identifies an exact same-law observational fiber.

\section{Setup}

Target (estimand / structural parameter): \(\mathcal D(P)=\{G':\exists(\mu',\alpha')\text{ such that }(G',\mu',\alpha',P)\in\mathcal M_d^\circ\}\).
Primary functional / estimator / diagnostic: \(P\mapsto t_4(P)\mapsto(A(P),F(P))\mapsto\{G_\rho:\rho(T)\in T,\ T\in\mathcal T_0(P)\}=\mathcal D(P)\).
\begin{itemize}
  \item \(d\) --- \texttt{integer}; \(\{2,3,\ldots\}\); \texttt{number of observed vertices}
  \item \(V\) --- \texttt{finite labeled vertex set}
  \par\smallskip\noindent\emph{Definition.} \(V=[d]\)
  \item \(G\) --- \texttt{directed acyclic graph}; the labeled DAGs on \(V\)
  \item \(E\) --- \texttt{directed edge set}
  \par\smallskip\noindent\emph{Definition.} \(E=E(G)\)
  \item \(\operatorname{Pa}_G\) --- \texttt{parent-\allowbreak{}set map}; \(V\to 2^V\)
  \par\smallskip\noindent\emph{Definition.} \(\operatorname{Pa}_G(j)=\{i:i\to j\in E\}\)
  \item \(\operatorname{Ch}_G\) --- \texttt{child-\allowbreak{}set map}; \(V\to 2^V\)
  \par\smallskip\noindent\emph{Definition.} \(\operatorname{Ch}_G(i)=\{j:i\to j\in E\}\)
  \item \(\alpha\) --- \texttt{edge-\allowbreak{}specific thinning vector}; \((0,1)^E\)
  \item \(\mu\) --- \texttt{innovation-\allowbreak{}rate vector}; \((0,\infty)^V\)
  \item \(B\) --- \texttt{edge-\allowbreak{}and-\allowbreak{}individual Bernoulli trial array}; \((i,j,r)\mapsto B_{ijr}\)
  \item \(\varepsilon\) --- \texttt{innovation random vector}; \(\mathbb N_0^V\)
  \item \(\circ\) --- \texttt{binomial thinning operator}; \((0,1)\times\mathbb N_0\to\mathbb N_0\)
  \par\smallskip\noindent\emph{Definition.} \(\alpha_{ij}\circ x=\sum_{r=1}^{x}B_{ijr}\)
  \item \(X\) --- \texttt{observed count random vector}; \(\mathbb N_0^V\)
  \item \(P\) --- \texttt{observational law}; \(\mathcal P(\mathbb N_0^V)\)
  \par\smallskip\noindent\emph{Definition.} \(P=\mathcal L(X)\)
  \item \(z\) --- \texttt{formal indeterminate vector}; \(\mathbb R^V\)
  \item \(H\) --- \texttt{root-\allowbreak{}cluster polynomial family}; \(V\to\mathbb R[z]\)
  \item \(\Psi_P\) --- \texttt{log probability-\allowbreak{}generating function}; \(z\mapsto\log E_P[\prod_{i\in V}z_i^{X_i}]\)
  \item \(c_k\) --- \texttt{exact-\allowbreak{}support log-\allowbreak{}PGF coefficient}; \(k\mapsto[z^k](\Psi_P(z)-\Psi_P(0))\)
  \item \(t_4(P)\) --- \texttt{degree-\allowbreak{}four coefficient signature}; \(\mathbb R^{\{k:1\le |k|\le4\}}\)
  \item \(\Gamma_i(k)\) --- finite set of rooted open-edge genealogies with label multiplicity \(k\)
  \item \(w_\gamma\) --- \texttt{genealogy probability weight}; \((0,1]\)
  \item \(\mathcal M_d^\circ\) --- \texttt{strict causally sufficient PB-\allowbreak{}SCM class}
  \item \(\mathcal D(P)\) --- \texttt{unrestricted observational DAG fiber}; a subset of the labeled DAGs on \(V\)
  \item \(S(P)\) --- \texttt{coefficient-\allowbreak{}recovered undirected skeleton}
  \item \(\mathsf{Head}_P\) --- \texttt{coefficient-\allowbreak{}recovered collider-\allowbreak{}head predicate}; \(V^3\to\{0,1\}\)
  \item \(A(P)\) --- \texttt{set of coefficient-\allowbreak{}compelled arrows entering multi-\allowbreak{}parent vertices}
  \item \(M(P)\) --- \texttt{set of recovered multi-\allowbreak{}parent vertices}
  \item \(F(P)\) --- \texttt{residual undirected graph}
  \item \(\mathcal T_0(P)\) --- \texttt{free residual-\allowbreak{}tree components}
  \item \(\rho\) --- \texttt{free-\allowbreak{}component root-\allowbreak{}choice map}; \(T\mapsto\rho(T)\in T\)
  \item \(G_\rho\) --- \texttt{root-\allowbreak{}choice completion}; the labeled DAGs on \(V\)
  \item \(\mathsf R_{u\to v}\) --- \texttt{source-\allowbreak{}edge reversal map}
  \item \(\theta_G(P)\) --- \texttt{fixed-\allowbreak{}DAG decoded parameter point}; \((0,\infty)^V\times(0,1)^{E(G)}\)
  \item \(n\) --- \texttt{sample size}; \(\mathbb N\)
  \item \(\eta\) --- \texttt{error probability}; \((0,1)\)
  \item \(X^{(b)}\) --- \texttt{sample observation}; \(\mathbb N_0^V\)
  \item \(I_4\) --- \texttt{low-\allowbreak{}count cell index set}
  \par\smallskip\noindent\emph{Definition.} \(I_4=\{x\in\mathbb N_0^V:|x|\le4\}\)
  \item \(p_x\) --- \texttt{count-\allowbreak{}cell probability}
  \par\smallskip\noindent\emph{Definition.} \(p_x=P(X=x)\)
  \item \(\widehat p_x\) --- \texttt{empirical count-\allowbreak{}cell probability}
  \par\smallskip\noindent\emph{Definition.} \(\widehat p_x=n^{-1}\sum_{b=1}^n\mathbf1\{X^{(b)}=x\}\)
  \item \(\epsilon_n\) --- \texttt{simultaneous cell radius}; \((0,\infty)\)
  \item \(Q_n\) --- \texttt{simultaneous truncated-\allowbreak{}cell confidence set}
  \item \(T(q)\) --- \texttt{rational cell-\allowbreak{}to-\allowbreak{}coefficient map}; \(q\mapsto(T_k(q))_{1\le|k|\le4}\)
  \item \(C_n\) --- \texttt{simultaneous degree-\allowbreak{}four coefficient region}
  \item \(\mathcal D(t)\) --- strict-model DAG fiber at coefficient vector \(t\)
  \item \(\mathcal D_n\) --- \texttt{outward-\allowbreak{}inverted graph set}
  \item \(\mathsf A_{\mathrm{out}}\) --- \texttt{certified outward-\allowbreak{}inversion handle}
\end{itemize}

\section{Assumptions}

\begin{assumption}[ass:structural-recursion]\label{ass:structural-recursion}
\(X_j=\varepsilon_j+\sum_{i\in\operatorname{Pa}_G(j)}\alpha_{ij}\circ X_i\quad(j\in V)\)
\par\smallskip\noindent\emph{Standard} (Poisson branching structural recursion, \cite{qiao2024causal}).
\end{assumption}

\begin{assumption}[ass:independent-branching-law]\label{ass:independent-branching-law}
\(\mathcal L\!\left(\varepsilon,(B_{ijr})_{(i,j)\in E,\,r\ge1}\right)=\bigotimes_{i\in V}\operatorname{Pois}(\mu_i)\otimes\bigotimes_{(i,j)\in E}\bigotimes_{r\ge1}\operatorname{Bernoulli}(\alpha_{ij})\)
\par\smallskip\noindent\emph{Standard} (Independent Poisson innovations and independent Bernoulli thinning trials, \cite{xiang2024identifiability}).
\end{assumption}

\begin{assumption}[ass:iid-sampling]\label{ass:iid-sampling}
\((X^{(b)})_{b=1}^n\stackrel{\mathrm{iid}}{\sim}P^{\otimes n}\)
\par\smallskip\noindent\emph{Standard} (Independent sampling from the observational count law, \cite{nakamuraPerezAbreu1993poisson}).
\end{assumption}

\section{Definitions}

\begin{definition}[def:strict-pb-class]\label{def:strict-pb-class}
\par\smallskip\noindent\emph{Defined object.} \texttt{Strict causally sufficient Poisson branching model class}
\par\smallskip\noindent\emph{Construction.} \(\mathcal M_d^\circ=\{(G,\mu,\alpha,P):G\text{ is a DAG on }V,\ \mu\in(0,\infty)^V,\ \alpha\in(0,1)^{E(G)},\ P=\mathcal L(X),\ \text{ass:structural-recursion and ass:independent-branching-law hold}\}\)
\par\smallskip\noindent\emph{Member properties.} \texttt{ass:\allowbreak{}structural-\allowbreak{}recursion}; \texttt{ass:\allowbreak{}independent-\allowbreak{}branching-\allowbreak{}law}.
\end{definition}

\begin{definition}[def:root-cluster]\label{def:root-cluster}
\par\smallskip\noindent\emph{Defined object.} \texttt{Root-\allowbreak{}cluster recursion}
\par\smallskip\noindent\emph{Construction.} \(H_i(z)=z_i\prod_{j\in\operatorname{Ch}_G(i)}(1-\alpha_{ij}+\alpha_{ij}H_j(z)),\qquad \Psi_P(z)=\sum_{i\in V}\mu_i(H_i(z)-1)\)
\par\smallskip\noindent\emph{Inputs.} \(G\); \(\mu\); \(\alpha\).
\end{definition}

\begin{definition}[def:log-pgf-coordinates]\label{def:log-pgf-coordinates}
\par\smallskip\noindent\emph{Defined object.} \texttt{Exact-\allowbreak{}support log-\allowbreak{}PGF coordinates}
\par\smallskip\noindent\emph{Construction.} \(c_k=[z^k](\Psi_P(z)-\Psi_P(0)),\qquad t_4(P)=(c_k)_{1\le |k|\le4}\)
\par\smallskip\noindent\emph{Inputs.} \(P\).
\end{definition}

\begin{definition}[def:genealogy-expansion]\label{def:genealogy-expansion}
\par\smallskip\noindent\emph{Defined object.} \texttt{Repeated-\allowbreak{}label genealogy weights}
\par\smallskip\noindent\emph{Construction.} \(\Gamma_i(k)=\{\gamma:\gamma\text{ is a finite open-edge descendant genealogy rooted at }i\text{ with observed-label multiplicity }k\},\quad w_\gamma=\prod_{e\in O_1(\gamma)}\alpha_e\prod_{e\in O_0(\gamma)}(1-\alpha_e)\)
\par\smallskip\noindent\emph{Inputs.} \(G\); \(\alpha\).
\end{definition}

\begin{definition}[def:observational-fiber]\label{def:observational-fiber}
\par\smallskip\noindent\emph{Defined object.} \texttt{Unrestricted observational DAG fiber}
\par\smallskip\noindent\emph{Construction.} \(\mathcal D(P)=\{G':\exists\mu'\in(0,\infty)^V,\ \exists\alpha'\in(0,1)^{E(G')}\text{ with }(G',\mu',\alpha',P)\in\mathcal M_d^\circ\}\)
\par\smallskip\noindent\emph{Inputs.} \(P\); \(\mathcal M_d^\circ\).
\end{definition}

\begin{definition}[def:signature-skeleton]\label{def:signature-skeleton}
\par\smallskip\noindent\emph{Defined object.} \texttt{Pair-\allowbreak{}coefficient skeleton}
\par\smallskip\noindent\emph{Construction.} \(S(P)=\{\{i,j\}\subseteq V:c_{e_i+e_j}>0\}\)
\par\smallskip\noindent\emph{Inputs.} \(t_4(P)\).
\end{definition}

\begin{definition}[def:signature-heads]\label{def:signature-heads}
\par\smallskip\noindent\emph{Defined object.} \texttt{Local collider-\allowbreak{}head predicate}
\par\smallskip\noindent\emph{Construction.} \(\mathsf{Head}_P(i,j;k)\Longleftrightarrow |\{i,j,k\}|=3,\ c_{e_i+e_k}c_{e_j+e_k}>0,\ \bigl[(c_{e_i+e_j}=0\wedge c_{e_i+e_j+e_k}=0)\vee(c_{e_i+e_j}>0\wedge c_{e_i+e_j+2e_k}>0)\bigr]\)
\par\smallskip\noindent\emph{Inputs.} \(t_4(P)\).
\end{definition}

\begin{definition}[def:signature-residual]\label{def:signature-residual}
\par\smallskip\noindent\emph{Defined object.} \texttt{Compelled incoming arrows and residual graph}
\par\smallskip\noindent\emph{Construction.} \(A(P)=\{i\to k:\exists j\in V,\ \mathsf{Head}_P(i,j;k)\},\quad M(P)=\{k:\exists i,j\in V,\ \mathsf{Head}_P(i,j;k)\},\quad F(P)=S(P)\setminus\{\{i,k\}:i\to k\in A(P)\}\)
\par\smallskip\noindent\emph{Inputs.} \(S(P)\); \(\mathsf{Head}_P\).
\end{definition}

\begin{definition}[def:root-choice-family]\label{def:root-choice-family}
\par\smallskip\noindent\emph{Defined object.} \texttt{Anchored and free root-\allowbreak{}choice completions}
\par\smallskip\noindent\emph{Construction.} \(\mathcal T_0(P)=\{T\in\operatorname{cc}(F(P)):T\cap M(P)=\varnothing\},\quad G_\rho=A(P)\cup\bigcup_{T\in\operatorname{cc}(F(P))}\operatorname{out}(T,r_T),\quad r_T=\begin{cases}m,&T\cap M(P)=\{m\},\\\rho(T),&T\in\mathcal T_0(P).\end{cases}\)
\par\smallskip\noindent\emph{Inputs.} \(A(P)\); \(F(P)\); \(M(P)\); \(\rho\).
\end{definition}

\begin{definition}[def:source-reversal]\label{def:source-reversal}
\par\smallskip\noindent\emph{Defined object.} \texttt{Positive source-\allowbreak{}edge reversal}
\par\smallskip\noindent\emph{Construction.} \(\mathsf R_{u\to v}(G,\mu,\alpha)=(G-u\to v+v\to u,\mu',\alpha'),\quad \mu'_u=\mu_u(1-\alpha_{uv}),\quad\mu'_v=\mu_v+\mu_u\alpha_{uv},\quad\alpha'_{vu}=\frac{\mu_u\alpha_{uv}}{\mu_v+\mu_u\alpha_{uv}},\quad(\mu'_i,\alpha'_e)=(\mu_i,\alpha_e)\text{ otherwise}\)
\par\smallskip\noindent\emph{Inputs.} \(G\); \(\mu\); \(\alpha\).
\end{definition}

\begin{definition}[def:rational-decoder]\label{def:rational-decoder}
\par\smallskip\noindent\emph{Defined object.} \texttt{Reverse-\allowbreak{}topological rational decoder}
\par\smallskip\noindent\emph{Construction.} \(\beta_i=\prod_{j\in\operatorname{Ch}_G(i)}(1-\alpha_{ij}),\qquad \alpha_{ij}=\frac{c_{e_i+e_j}}{c_{e_i+e_j}+c_{e_i}\beta_j},\qquad \mu_i=\frac{c_{e_i}}{\beta_i}\),\text{ evaluated in reverse topological order}\)
\par\smallskip\noindent\emph{Inputs.} \(G\); \(c_k\).
\end{definition}

\begin{definition}[def:triangle-witness]\label{def:triangle-witness}
\par\smallskip\noindent\emph{Defined object.} \texttt{Strict degree-\allowbreak{}three collision witness}
\par\smallskip\noindent\emph{Construction.} \(\mu_s=6,\ \mu_m=2,\ \mu_r=1,\ \alpha_{sm}=2/3,\ \alpha_{sr}=1/2,\ \alpha_{mr}=1/2\) on \(s\to m\to r\) and \(s\to r\), for each choice of sink \(r\) among three labels
\par\smallskip\noindent\emph{Inputs.} \(V\).
\end{definition}

\begin{definition}[def:count-cell-map]\label{def:count-cell-map}
\par\smallskip\noindent\emph{Defined object.} \texttt{Finite rational count-\allowbreak{}cell map}
\par\smallskip\noindent\emph{Construction.} \(I_4=\{x\in\mathbb N_0^V:|x|\le4\},\quad T_k(q)=\sum_{\ell=1}^{|k|}\frac{(-1)^{\ell+1}}{\ell}\sum_{x_1+\cdots+x_\ell=k,\,x_a\ne0}\prod_{a=1}^{\ell}\frac{q_{x_a}}{q_0},\quad1\le|k|\le4\)
\par\smallskip\noindent\emph{Inputs.} \(q_x:x\in I_4\).
\end{definition}

\begin{definition}[def:coefficient-region]\label{def:coefficient-region}
\par\smallskip\noindent\emph{Defined object.} \texttt{Simultaneous coefficient confidence region}
\par\smallskip\noindent\emph{Construction.} \(\epsilon_n=\sqrt{\log(2|I_4|/\eta)/(2n)},\quad Q_n=\{q:q_x\ge0,\ \sum_{x\in I_4}q_x\le1,\ q_0>0,\ |q_x-\widehat p_x|\le\epsilon_n\ (x\in I_4)\},\quad C_n=\{T(q):q\in Q_n\}\)
\par\smallskip\noindent\emph{Inputs.} \(X^{(b)}\); \(n\); \(\eta\).
\end{definition}

\begin{definition}[def:outer-fiber]\label{def:outer-fiber}
\par\smallskip\noindent\emph{Defined object.} \texttt{Outward-\allowbreak{}inverted graph set}
\par\smallskip\noindent\emph{Construction.} \(\mathcal D(t)=\{G:\exists(G,\mu,\alpha,P')\in\mathcal M_d^\circ,\ t_4(P')=t\},\qquad \mathcal D_n=\bigcup_{t\in C_n}\mathcal D(t)\)
\par\smallskip\noindent\emph{Inputs.} \(C_n\); \(\mathcal M_d^\circ\).
\end{definition}

\begin{definition}[def:outward-inversion-handle]\label{def:outward-inversion-handle}
\par\smallskip\noindent\emph{Defined object.} \texttt{Certified source-\allowbreak{}forest outward-\allowbreak{}inversion handle}
\par\smallskip\noindent\emph{Construction.} \(\mathsf A_{\mathrm{out}}(n,\eta,\widehat p;O,\kappa)\) is a specification for one fixed rational confidence instance: \(n\in\mathbb N\), \(\eta\in\mathbb Q\), a rational cell-frequency map \(\widehat p=(\widehat p_x)_x\), an output list \(O\) of labeled edge functions, and a certificate family \(\kappa=(\kappa_e)_e\) indexed by every labeled edge function \(e\). It holds exactly when: (i) \(e\in O\Longleftrightarrow e\in\mathcal D_n(n,\eta,\widehat p)\) for every \(e\), with duplicate list entries immaterial; and (ii) each \(\kappa_e\) carries a Boolean verdict, singleton/pair and degree-three/four coefficient intervals, a source-forest candidate list, a reverse-topological order, decoded rate and edge-probability intervals, and a branch-and-bound trace satisfying all of the following correctness clauses: the interval lists cover their advertised degrees and bound every \(t\in C_n(n,\eta,\widehat p)\); the source-forest list equals \(\mathcal D_n(n,\eta,\widehat p)\) and every listed graph has a strict-law root-choice witness with signature in \(C_n\); a true verdict certifies acyclicity, a valid reverse-topological order, and decoded intervals for every vertex and present edge, with those intervals correct for every \(t\in C_n\) realizing \(e\); the branch-and-bound trace starts from the source-forest list, decreases by inclusion, preserves \(\mathcal D_n\) throughout, and ends at \(O\); and the certificate verdict is true exactly when \(e\in\mathcal D_n\), equivalently true exactly when \(e\in O\) and false exactly when \(e\notin O\). This definition is only an interface predicate and asserts neither existence nor polynomial-delay construction.
\par\smallskip\noindent\emph{Inputs.} \(n\); \(\eta\); \(\widehat p_x\).
\end{definition}

\section{Main results}

\begin{lemma}[lem:genealogy-coefficient]\label{lem:genealogy-coefficient}
For every \((G,\mu,\alpha,P)\in\mathcal M_d^\circ\) and every nonzero multi-index \(k\), the exact-support coefficient satisfies \(c_k=\sum_{i\in V}\mu_i\sum_{\gamma\in\Gamma_i(k)}w_\gamma\). The identity counts genealogical occurrences before observed labels are merged, so distinct downstream paths that repeat a vertex label contribute to the same exponent without cancellation.
\end{lemma}
\begin{proof}
Fix \((G,\mu,\alpha,P)\in\mathcal M_d^\circ\) and a nonzero multi-index \(k\).  We make the occurrence bookkeeping explicit.  Start from one individual labeled \(i\).  Whenever an occurrence has label \(a\), perform, independently, one Bernoulli trial for each edge \(a\to b\); an open trial creates one new occurrence labeled \(b\), and a closed trial creates none.  Acyclicity implies that every lineage has length at most \(d-1\).  Because \(V\) and every outdegree are finite, each resulting genealogy and the collection of its possible open--closed histories are finite.

Let \(K_i(\omega)\in\mathbb N_0^V\) count the labels of the occurrences in a history \(\omega\) rooted at \(i\).  We prove by reverse topological induction that
\[
 H_i(z)=\mathbb E\bigl[z^{K_i(\omega)}\bigr].
\]
At a sink, both sides equal \(z_i\).  At a general vertex, the root occurrence contributes \(z_i\).  For each \(j\in\operatorname{Ch}_G(i)\), the edge trial is closed with probability \(1-\alpha_{ij}\), while, if it is open, the child occurrence initiates a cluster with probability-generating polynomial \(H_j\) by the induction hypothesis.  Independence of all occurrence-indexed trials in ass:independent-branching-law therefore gives
\[
 \mathbb E\bigl[z^{K_i(\omega)}\bigr]
 =z_i\prod_{j\in\operatorname{Ch}_G(i)}
       \bigl(1-\alpha_{ij}+\alpha_{ij}H_j(z)\bigr)
 =H_i(z),
\]
where the last equality is def:root-cluster.  In this induction occurrences are never identified merely because they have the same observed label.  Thus two open directed paths ending at the same vertex produce two occurrences and contribute two to the corresponding coordinate of \(K_i(\omega)\).

Interpret \(O_1(\omega)\) and \(O_0(\omega)\) in def:genealogy-expansion as occurrence-indexed collections: two trials of the same graphical edge made by distinct occurrences appear as distinct factors.  Ass:independent-branching-law then yields
\[
 \Pr(\omega)
 =\prod_{e\in O_1(\omega)}\alpha_e
   \prod_{e\in O_0(\omega)}(1-\alpha_e)
 =w_\omega.
\]
Hence finite coefficient extraction gives
\[
 [z^k]H_i(z)
 =\sum_{\omega:K_i(\omega)=k}\Pr(\omega)
 =\sum_{\gamma\in\Gamma_i(k)}w_\gamma.
\tag{1}
\]

It remains to aggregate innovations.  Conditional on \(\varepsilon_i=m\), the \(m\) innovation individuals at \(i\) initiate independent copies of this cluster, so their joint probability-generating polynomial is \(H_i(z)^m\).  Independence across \(i\), followed by the Poisson probability-generating identity, gives, for \(z\in[0,1]^V\),
\[
 \mathbb E_P[z^X]
 =\prod_{i\in V}\mathbb E\bigl[H_i(z)^{\varepsilon_i}\bigr]
 =\prod_{i\in V}\exp\{\mu_i(H_i(z)-1)\}.
\]
Taking logarithms proves the log-PGF identity in def:root-cluster,
\[
 \Psi_P(z)=\sum_{i\in V}\mu_i(H_i(z)-1).
\tag{2}
\]
Both sides of (2) are polynomials, so the identity on \([0,1]^V\) is also a polynomial identity.  Since \(k\ne0\), subtracting the constant \(\Psi_P(0)\) does not alter its \(z^k\)-coefficient.  Equations (1)--(2) and def:log-pgf-coordinates therefore imply
\[
 c_k
 =\sum_{i\in V}\mu_i[z^k]H_i(z)
 =\sum_{i\in V}\mu_i\sum_{\gamma\in\Gamma_i(k)}w_\gamma.
\]
Finally, strictness in def:strict-pb-class gives \(\mu_i>0\), \(0<\alpha_e<1\), and hence \(w_\gamma>0\) for every finite genealogy.  Thus repeated-label contributions add without cancellation, as claimed.
\end{proof}
\par\smallskip\noindent \textit{Justification.} This is the load-bearing bridge from arbitrary-DAG branching genealogies to every local coefficient certificate. \textit{Gap.} Xiang et al. give a path-to-monomial correspondence, but the repeated-label multiplicity bookkeeping needed for a universal low-degree converse is not isolated in their Proposition 1. \textit{Consumer.} The source-forest converse and the degree-four compression theorem both use this formula to rule out false signatures from overlapping descendants.

\begin{theorem}[thm:local-signatures]\label{thm:local-signatures}
For every strict-model law \(P\), \(c_{e_i+e_j}>0\) if and only if \(\{i,j\}\in S(P)\). On every induced path \(i-k-j\), \(c_{e_i+e_j+e_k}=0\) if and only if \(i\to k\leftarrow j\). On every triangle, \(c_{e_i+e_j+2e_k}>0\) if and only if \(i\to k\leftarrow j\). Consequently \(A(P)\) is exactly the union of all arrows entering every vertex of indegree at least two in every member of \(\mathcal D(P)\).
\end{theorem}
\begin{proof}
Fix a strict representation \((G,\mu,\alpha,P)\).  By lem:genealogy-coefficient and strictness, \(c_k>0\) holds exactly when a one-innovation genealogy with label-multiplicity vector \(k\) exists.  Its root label occurs, and every label on every open directed path in it occurs.  If \(i=j\), acyclicity prevents a genealogy rooted at \(i\) from producing a second occurrence of \(i\), while a genealogy rooted elsewhere also contains its distinct root label; hence \(c_{2e_i}=0\), consistently with the absence of a self-edge.  It remains to consider distinct \(i,j\).

A genealogy with multiplicity \(e_i+e_j\) must be rooted at \(i\) or \(j\).  A directed path from the root to the other label cannot have an intermediate vertex, since that label would also occur.  Therefore \(i\) and \(j\) must be adjacent.  Conversely, if \(i\to j\), open that edge and close every other outgoing trial from the occurrences at \(i\) and \(j\).  There are finitely many such trials and each required outcome has positive probability because \(0<\alpha_e<1\); the resulting multiplicity is \(e_i+e_j\).  The case \(j\to i\) is symmetric.  Hence
\[
 c_{e_i+e_j}>0
 \quad\Longleftrightarrow\quad
 \{i,j\}\text{ is an edge of the undirected skeleton of }G.
\tag{3}
\]
By def:signature-skeleton, \(S(P)\) is therefore that skeleton.  Because the left side of (3) is a functional of \(P\), this skeleton is common to every \(G\in\mathcal D(P)\).

Now let \(i-k-j\) be an induced path in this skeleton.  Exact support on \(\{i,j,k\}\) forces the genealogy root and every traversed vertex to lie in this triple, so only the two displayed skeleton edges can be used.  If \(i\to k\leftarrow j\), no vertex in the triple can reach both other vertices, and no genealogy contains all three labels.  Thus \(c_{e_i+e_j+e_k}=0\).  Each of the other three acyclic orientations is an out-tree: it is rooted at \(i\), \(j\), or \(k\).  Opening its two tree edges and closing every other outgoing trial of the three occurrences produces each label once with positive probability.  Consequently
\[
 c_{e_i+e_j+e_k}=0
 \quad\Longleftrightarrow\quad
 i\to k\leftarrow j.
\tag{4}
\]

Next suppose that \(i,j,k\) form a triangle.  Every acyclic orientation of a triangle is transitive; write its unique order as \(s\to m\to r\), together with \(s\to r\).  If \(k=r\), opening all three triangle edges in the cluster rooted at \(s\), and closing all other outgoing trials of all created occurrences, gives multiplicity \(e_s+e_m+2e_r\) with positive probability.  Conversely, exact support again confines the root and all paths to the triangle.  In its transitive orientation, the only vertex reachable from one root along two distinct directed paths is the sink \(r\), reached from \(s\) both directly and through \(m\).  Since a Bernoulli edge trial creates at most one child occurrence, no other label can occur twice in a genealogy having the other two labels once.  Therefore
\[
 c_{e_i+e_j+2e_k}>0
 \quad\Longleftrightarrow\quad
 i\to k\leftarrow j.
\tag{5}
\]
The occurrence-level interpretation from lem:genealogy-coefficient is essential here: the two paths to the sink yield two occurrences even when their observed labels agree, and descendants outside the triangle cannot create a false positive because their labels would violate exact support.

Equations (3)--(5) and def:signature-heads show that \(\mathsf{Head}_P(i,j;k)\) holds exactly when \(i\) and \(j\) are distinct parents of \(k\).  Indeed, (3) makes the two pair factors positive precisely when both skeleton adjacencies are present; if \(i,j\) are nonadjacent, the first disjunct in def:signature-heads is equivalent to (4), while if they are adjacent, its second disjunct is equivalent to (5).  Thus, for every vertex \(k\) of indegree at least two and every parent \(i\) of \(k\), pairing \(i\) with any other parent places \(i\to k\) in \(A(P)\).  Conversely, every arrow placed in \(A(P)\) has a head certificate of this form and therefore enters a vertex with at least two parents.  All coefficients used are law functionals, so the conclusion holds simultaneously in every member of \(\mathcal D(P)\).  No conditional-independence or faithfulness premise is used.
\end{proof}
\par\smallskip\noindent \textit{Justification.} The theorem converts degree-four observable coordinates into the complete set of fixed incoming arrows. \textit{Gap.} Xiang et al. establish the local tests, while the claim here assembles them over every parent pair without faithfulness and makes their fiber-wide invariance explicit. \textit{Consumer.} The residual-forest theorem uses \(A(P)\) as its observable anchor set.

\begin{theorem}[thm:source-reversal]\label{thm:source-reversal}
If \(u\) is a source of \(G\), \(\operatorname{Pa}_G(v)=\{u\}\), and \(u\to v\in E\), then \(\mathsf R_{u\to v}\) is acyclic, all transformed rates and thinning probabilities are strict, and its complete log-PGF equals \(\Psi_P\) identically, including the constant coefficient. The equality remains valid when \(u\) and \(v\) have arbitrarily overlapping descendant sets.
\end{theorem}
\begin{proof}
Put \(t=\alpha_{uv}\).  If there were a directed path from \(u\) to \(v\) other than the edge \(u\to v\), its last edge would enter \(v\) from a parent different from \(u\), contradicting \(\operatorname{Pa}_G(v)=\{u\}\).  Replacing \(u\to v\) by \(v\to u\) therefore cannot create a directed cycle.  Moreover, after reversal \(v\) is a source and \(u\) has sole parent \(v\).

Define the following formal polynomials using the unchanged descendant root-cluster polynomials:
\[
 U:=z_u\prod_{a\in\operatorname{Ch}_G(u)\setminus\{v\}}
       (1-\alpha_{ua}+\alpha_{ua}H_a),
 \qquad
 W:=z_v\prod_{b\in\operatorname{Ch}_G(v)}
       (1-\alpha_{vb}+\alpha_{vb}H_b).
\]
Acyclicity implies that neither product contains \(u\) or \(v\) as a descendant.  The products may involve arbitrarily many common descendant variables; no disjointness is being assumed.  The old root polynomials are
\[
 H_u=U(1-t+tW),\qquad H_v=W.
\]
With
\[
 \mu'_u=\mu_u(1-t),\qquad
 \mu'_v=\mu_v+\mu_ut,
 \qquad
 s:=\alpha'_{vu}=\frac{\mu_ut}{\mu_v+\mu_ut},
\]
the new root polynomials are
\[
 H'_u=U,\qquad H'_v=W(1-s+sU).
\]
No third vertex is an ancestor of \(u\), because \(u\) is a source, and no third vertex is an ancestor of \(v\), because its sole parent is \(u\).  The same assertion holds after reversal with \(u,v\) interchanged.  Therefore every \(H_a\) with \(a\notin\{u,v\}\) is unchanged.

The two nonconstant contributions satisfy the formal-polynomial identity
\[
\begin{aligned}
 \mu_uH_u+\mu_vH_v
 &=\mu_uU(1-t+tW)+\mu_vW\\
 &=\mu_u(1-t)U+\mu_vW+\mu_utUW\\
 &=\mu'_uU+\mu'_vW(1-s+sU)\\
 &=\mu'_uH'_u+\mu'_vH'_v,
\end{aligned}
\]
where \(\mu'_v(1-s)=\mu_v\) and \(\mu'_vs=\mu_ut\).  Also
\[
 \mu'_u+\mu'_v=\mu_u+\mu_v.
\]
Substituting these two identities into
\(\Psi(z)=\sum_a\mu_a(H_a(z)-1)\) proves equality of the complete log-PGFs, including the constant coefficient.

Finally, \(\mu'_u>0\), \(\mu'_v>0\), and
\[
 0<\frac{\mu_ut}{\mu_v+\mu_ut}<1
\]
because \(\mu_u,\mu_v>0\) and \(0<t<1\); every unchanged parameter remains strict.  This proves all assertions, and the algebra explicitly permits overlapping descendant sets.
\end{proof}
\par\smallskip\noindent \textit{Justification.} This transformation supplies a positive same-law witness for every admissible ambiguity. \textit{Gap.} Qiao et al. give the two-variable reverse representation; they do not prove a complete-PGF identity inside an arbitrary graph with shared descendants. \textit{Consumer.} Every free-tree rerooting in the complete fiber obtains an explicit alternative mechanism and parameter point.

\begin{theorem}[thm:fixed-dag-decoder]\label{thm:fixed-dag-decoder}
For any \(G\in\mathcal D(P)\), the reverse-topological recursion in \(\theta_G(P)\) returns the unique \((\mu,\alpha)\in(0,\infty)^V\times(0,1)^{E(G)}\) inducing \(P\). Every denominator is positive, and the fixed-DAG coefficient map has full column rank throughout the strict parameter domain.
\end{theorem}
\begin{proof}
Let \(G\in\mathcal D(P)\).  By def:observational-fiber, at least one strict parameter point \((\mu,\alpha)\) on \(G\) induces \(P\).  Put
\[
 \beta_i:=\prod_{j\in\operatorname{Ch}_G(i)}(1-\alpha_{ij}).
\]
The coefficient of \(z_i\) in def:root-cluster can only arise from an innovation at \(i\) for which all outgoing trials are closed; a cluster rooted anywhere else contains its distinct root label.  Hence
\[
 c_{e_i}=\mu_i\beta_i>0.
\tag{6}
\]
For an edge \(i\to j\), exact support \(\{i,j\}\) requires the trial \(i\to j\) to be open and all other outgoing trials of the occurrences at \(i\) and \(j\) to be closed.  Acyclicity excludes a cluster rooted at \(j\) that also contains \(i\), and a root outside the pair violates exact support.  Thus
\[
 \begin{aligned}
 c_{e_i+e_j}
 &=\mu_i\alpha_{ij}
   \prod_{h\in\operatorname{Ch}_G(i)\setminus\{j\}}(1-\alpha_{ih})\,\beta_j\\
 &=\frac{c_{e_i}\alpha_{ij}\beta_j}{1-\alpha_{ij}}>0.
 \end{aligned}
\tag{7}
\]
Solving (7), with \(1-\alpha_{ij}>0\), yields
\[
 \alpha_{ij}
 =\frac{c_{e_i+e_j}}
        {c_{e_i+e_j}+c_{e_i}\beta_j}.
\tag{8}
\]

Evaluate (8) in reverse topological order.  At every sink, \(\beta_i=1\).  Once all children of \(i\) have been processed, each \(\beta_j\) on the right side of (8) is known, so all outgoing \(\alpha_{ij}\) are determined; their product determines \(\beta_i\), after which (6) gives
\[
 \mu_i=\frac{c_{e_i}}{\beta_i}.
\tag{9}
\]
Every denominator in (8) is positive by (6)--(7); its numerator is positive and its denominator is strictly larger than its numerator, so the decoded \(\alpha_{ij}\) lies in \((0,1)\).  Every \(\beta_i\) and every decoded \(\mu_i\) is positive.  Induction through the reverse topological order proves that def:rational-decoder returns the generating parameter point.  Applying the same recursion to any two strict parameter points on \(G\) inducing \(P\) proves edge-by-edge and then vertex-by-vertex equality, hence uniqueness.

For the rank assertion, write \(m=|E(G)|\), let \(\Theta_G=(0,\infty)^V\times(0,1)^{E(G)}\), and let \(\Phi_G\) denote the fixed-DAG coefficient map.  Let \(\pi\Phi_G\) retain only the \(d\) singleton coordinates and the \(m\) pair coordinates corresponding to edges of \(G\).  The recursion (8)--(9) defines a rational map \(R_G\) on an open neighborhood of every point of \(\pi\Phi_G(\Theta_G)\), because all displayed denominators are strictly positive there, and the preceding calculation proves
\[
 R_G\circ\pi\Phi_G=\operatorname{id}_{\Theta_G}.
\]
Both maps are differentiable on the strict domain.  The chain rule gives
\[
 D R_G\,D(\pi\Phi_G)=I_{d+m}.
\]
Therefore \(D(\pi\Phi_G)\), and hence the Jacobian of the full coefficient map \(D\Phi_G\), has column rank \(d+m\) everywhere on \(\Theta_G\).
\end{proof}
\par\smallskip\noindent \textit{Justification.} Graph ambiguity is separated cleanly from numerical nonidentification by a constructive inverse. \textit{Gap.} Xiang et al. estimate PGF statistics for graph learning but do not state a global fixed-DAG rational inverse or parameter-fiber injectivity theorem. \textit{Consumer.} Analysts can decode and compare every mechanism in the football-event and shopping-mall observational fibers.

\begin{theorem}[thm:source-forest-fiber]\label{thm:source-forest-fiber}
For every \(P\) generated by \(\mathcal M_d^\circ\), \(F(P)\) is a forest, each connected component of \(F(P)\) contains at most one vertex of \(M(P)\), and \(\mathcal D(P)=\{G_\rho:\rho(T)\in T\text{ for every }T\in\mathcal T_0(P)\}\). A component meeting \(M(P)\) is directed away from its unique anchor, and each free component is independently directed away from its chosen root. This necessary-and-sufficient classification strictly strengthens the sufficient local orientation conclusions of Xiang et al. and Qiao et al.
\end{theorem}
\begin{proof}
Fix a strict-model law \(P\) and any \(G\in\mathcal D(P)\).  By thm:local-signatures, \(S(P)\) is the undirected skeleton recovered by the pair coefficients and \(A(P)\) consists exactly of the arrows whose heads have indegree at least two.  Delete those arrows while retaining the orientations of the remaining edges of \(F(P)\).  Every vertex in \(M(P)\) then has residual indegree zero, and every vertex outside \(M(P)\) has residual indegree at most one.

If the undirected residual graph contained a cycle, choose the latest vertex on that cycle in a topological ordering of \(G\).  Both cycle edges incident to it would point inward, giving residual indegree at least two, a contradiction.  Thus \(F(P)\) is a forest.

Let \(T\) be a connected component of \(F(P)\), and put \(r=|T|\).  Since \(T\) is a tree, it has \(r-1\) edges, whence
\[
 \sum_{v\in T}\operatorname{indeg}_T(v)=r-1.
\tag{10}
\]
Every summand is at most one.  Equation (10) therefore implies that exactly one vertex has residual indegree zero and all others have residual indegree one.  Every member of \(T\cap M(P)\) has residual indegree zero, so \(|T\cap M(P)|\le1\).  If \(T\cap M(P)=\{m\}\), then \(m\) is the unique residual root.  The first edge on every path from \(m\) must point away from \(m\); after a vertex on the path has received its unique incoming residual edge, the next edge must point away from it.  Induction along the unique tree paths directs all of \(T\) away from \(m\).  If \(T\cap M(P)=\varnothing\), the same argument directs \(T\) away from its unique residual root, which may a priori be any vertex of \(T\).

No arrow enters a free component \(T\) from outside.  A deleted arrow has its head in \(M(P)\), whereas \(T\cap M(P)=\varnothing\); an undeleted edge incident to \(T\) belongs to \(F(P)\) and would put its other endpoint in the same connected component.  Hence the residual root of a free component is a global source of \(G\), while every other vertex of that component has its tree predecessor as its sole parent.  Since the signature objects are functionals of \(P\), these conclusions hold for every \(G\in\mathcal D(P)\), proving
\[
 \mathcal D(P)\subseteq
 \{G_\rho:\rho(T)\in T\text{ for all }T\in\mathcal T_0(P)\}.
\tag{11}
\]

For the reverse inclusion, begin with the fixed representation \(G\).  In a free tree having current root \(v_0\), let \(v_0,v_1,\ldots,v_\ell\) be the unique path to a desired root \(v_\ell\).  Initially \(v_0\) is a global source and \(\operatorname{Pa}_G(v_1)=\{v_0\}\), so thm:source-reversal applies to \(v_0\to v_1\).  After this reversal, \(v_1\) is the global source of the component and \(v_2\), if present, still has sole parent \(v_1\).  Inductively, each map
\[
 \mathsf R_{v_0\to v_1},\quad
 \mathsf R_{v_1\to v_2},\quad\ldots,\quad
 \mathsf R_{v_{\ell-1}\to v_\ell}
\tag{12}
\]
is applicable.  By thm:source-reversal, every step is acyclic, preserves strict positivity, and preserves the complete log-PGF even with overlapping descendant sets.  It changes no edge outside that free tree.  Distinct free components have disjoint residual edge sets and no incoming edges, so the construction (12) may be performed independently in all of them.  The resulting graph is precisely \(G_\rho\) and has a strict parameter point inducing \(P\).  This proves the reverse inclusion in (11), and hence
\[
 \mathcal D(P)=
 \{G_\rho:\rho(T)\in T\text{ for every }T\in\mathcal T_0(P)\}.
\]
This is a necessary-and-sufficient global classification, whereas the cited pathwise and local orientation conclusions of Xiang et al.\ \cite[Theorem~7]{xiang2024identifiability} and Qiao et al.\ \cite[Theorem~7]{qiao2024causal} are sufficient orientation rules; the exhaustive converse and the independently rerootable free components are the strict additional conclusions.
\end{proof}
\par\smallskip\noindent \textit{Justification.} This is the theorem-level kernel: the complete one-law observational equivalence class is an explicit source-forest root-choice family. \textit{Gap.} Qiao et al. and Xiang et al. give sufficient path or local PGF orientations, but neither paper proves an arbitrary-DAG converse or exhausts every compatible representation. \textit{Consumer.} The result turns the remaining unoriented football and shopping arrows into certified ambiguity classes instead of unresolved algorithm output.

\begin{proposition}[prop:fiber-cardinality-certificates]\label{prop:fiber-cardinality-certificates}
The exact fiber size is \(|\mathcal D(P)|=\prod_{T\in\mathcal T_0(P)}|T|\). An edge is compelled if and only if it does not lie in a free tree; its certificate is either a local head signature or the unique path from an anchor. Every free-tree edge admits both directions, certified by a finite sequence of maps \(\mathsf R_{u\to v}\).
\end{proposition}
\begin{proof}
By thm:source-forest-fiber, a compatible DAG is specified by choosing exactly one root in each free residual tree, while every anchored component is fixed.  Root choices in distinct components affect disjoint edge sets.  Two distinct roots in one tree give distinct orientations because the root is the unique vertex of residual indegree zero.  Thus the root-choice map is both injective and surjective, and finiteness gives
\[
 |\mathcal D(P)|=\prod_{T\in\mathcal T_0(P)}|T|,
\]
where the empty product equals one.

Every arrow in \(A(P)\) is compelled by the corresponding local head signature from thm:local-signatures.  In an anchored residual tree, that signature fixes the unique anchor, and each residual edge is forced away from the anchor along the unique tree path; the head signature together with this path is a finite certificate of the edge orientation.  Hence every edge not lying in a free tree is compelled.

Conversely, let \(\{a,b\}\) be an edge of a free tree.  Choosing \(a\) as the root orients it as \(a\to b\), while choosing \(b\) as the root orients it as \(b\to a\).  Both completions belong to \(\mathcal D(P)\) by thm:source-forest-fiber, so neither orientation is compelled.  Constructively, move any current root along the unique path to \(a\), and then, when the opposite direction is desired, across \(a-b\) to \(b\).  At each step the current root is a global source and the next vertex has it as sole parent; hence thm:source-reversal certifies the step by a map \(\mathsf R_{u\to v}\) that preserves the complete law and strict parameters.  The resulting finite reversal sequence certifies both directions and proves the claimed if-and-only-if characterization.
\end{proof}
\par\smallskip\noindent \textit{Justification.} Cardinality and certificates make the equivalence theorem directly auditable at the edge level. \textit{Gap.} No cited PB-SCM paper counts the full fiber or supplies constructive certificates for both compelled and reversible arrows. \textit{Consumer.} Applied users can report which event pathways are invariant across all mechanisms and enumerate every admissible alternative.

\begin{theorem}[thm:degree-four-completeness]\label{thm:degree-four-completeness}
For any two laws \(P,P'\) generated by \(\mathcal M_d^\circ\), \(t_4(P)=t_4(P')\) if and only if \(P=P'\). Equality of \(t_4\) also implies \(\mathcal D(P)=\mathcal D(P')\) and \(\theta_G(P)=\theta_G(P')\) for every compatible \(G\).
\end{theorem}
\begin{proof}
If \(P=P'\), then their log-PGFs coincide, so def:log-pgf-coordinates immediately gives \(t_4(P)=t_4(P')\).  Conversely, assume \(t_4(P)=t_4(P')\).  The constructions in def:signature-skeleton, def:signature-heads, and def:signature-residual use only coefficients of total degree at most four.  Thus the two laws have the same \(S\), \(\mathsf{Head}\), \(A\), \(M\), \(F\), and collection \(\mathcal T_0\).  Applying thm:source-forest-fiber to each law yields the same nonempty root-choice family,
\[
 \mathcal D(P)=\{G_\rho\}=\mathcal D(P').
\tag{13}
\]
Choose any \(G\) in this common family; nonemptiness follows because both laws are generated by \(\mathcal M_d^\circ\).  The rational recursion of thm:fixed-dag-decoder uses only singleton and skeleton-edge pair coefficients, all contained in the common \(t_4\).  Its uniqueness conclusion therefore gives
\[
 \theta_G(P)=\theta_G(P').
\tag{14}
\]
In particular, the two parameter points have identical \(\mu\) and \(\alpha\).  Substituting (14) into def:root-cluster gives equality of the complete log-PGFs, including the constant term:
\[
 \Psi_P(z)=\sum_{i\in V}\mu_i(H_i(z)-1)=\Psi_{P'}(z).
\]
Exponentiation gives equality of the probability-generating functions on \([0,1]^V\).  Since
\[
 \mathbb E_P[z^X]=\sum_{x\in\mathbb N_0^V}P(X=x)z^x
\]
and the multivariate power series is absolutely convergent on \([0,1)^V\), uniqueness of its coefficients implies \(P(X=x)=P'(X=x)\) for every \(x\), hence \(P=P'\).  Equations (13)--(14) also prove the asserted equality of fibers and of decoded parameter points for every compatible \(G\).
\end{proof}
\par\smallskip\noindent \textit{Justification.} A fixed low-order signature determines an arbitrarily high-degree branching law and its entire graph-parameter fiber. \textit{Gap.} Qiao et al. use a path-dependent cumulant order and Xiang et al. use the complete PGF; neither gives a universal finite cutoff. \textit{Consumer.} The full population identification problem reduces to finitely many count-cell functionals for any fixed number of vertices.

\begin{theorem}[thm:degree-three-sharpness]\label{thm:degree-three-sharpness}
The three strict triangle laws in the construction \(\text{def:triangle-witness}\) have identical \(\Psi_P(0)\) and identical \(c_k\) for every \(|k|\le3\), but their coefficients of \(z_s z_m z_r^2\) differ and their DAG fibers have different triangle sinks. Adding independent positive-Poisson coordinates extends the failure to every \(d\ge3\).
\end{theorem}
\begin{proof}
Let \(L=\{1,2,3\}\), fix a proposed sink \(r\in L\), and write \(L\setminus\{r\}=\{s,m\}\), with either ordering of the two nonsink labels.  The graph in \texttt{def:triangle-witness} is the acyclic triangle
\[
s\longrightarrow m\longrightarrow r,
\qquad s\longrightarrow r.
\]
All three innovation rates are positive and all three thinning probabilities belong to \((0,1)\), so, with the independent innovations and trials required by \texttt{def:strict-pb-class}, this is a strict model.

We first calculate its law.  By the recursion in \texttt{def:root-cluster},
\[
H_r=z_r,
\qquad
H_m=z_m\left(\frac12+\frac12z_r\right),
\]
and
\[
\begin{aligned}
H_s
&=z_s\left(\frac13+\frac23H_m\right)
          \left(\frac12+\frac12z_r\right)\\
&=z_s\left(\frac13+\frac13z_m+\frac13z_mz_r\right)
          \left(\frac12+\frac12z_r\right)\\
&=\frac16z_s+\frac16z_sz_m+\frac16z_sz_r
  +\frac13z_sz_mz_r+\frac16z_sz_mz_r^2.
\end{aligned}
\]
Consequently, again by \texttt{def:root-cluster},
\[
\begin{aligned}
\Psi^{(r)}(z)
&=6(H_s-1)+2(H_m-1)+(H_r-1)\\
&=-9+\sum_{v\in L}z_v
  +\sum_{\{u,v\}\subset L}z_uz_v
  +2\prod_{v\in L}z_v
  +\left(\prod_{v\in L}z_v\right)z_r.
\end{aligned}
\tag{1}
\]
In particular, \(\Psi^{(r)}(0)=-9\).  The portion of (1) of total degree at most three is symmetric in the three labels and therefore does not depend on \(r\).  By \texttt{def:log-pgf-coordinates}, this proves that the three laws have the same \(c_k\) for every multi-index \(k\) with \(|k|\leq 3\).  This includes the zero coefficient after subtraction of \(\Psi^{(r)}(0)\).  On the other hand, for every fixed \(q\in L\),
\[
\left[\left(\prod_{v\in L}z_v\right)z_q\right]
\bigl(\Psi^{(r)}(z)-\Psi^{(r)}(0)\bigr)
=\mathbf 1\{q=r\}.
\tag{2}
\]
Thus, if \(r\) denotes the sink of one of the three laws, its coefficient of \(z_sz_mz_r^2\) is one, whereas that coefficient is zero for each law whose sink is one of the other two fixed labels.  Hence the three full laws are distinct although all their degree-at-most-three coordinates agree.

For completeness, we now prove directly that the differing fourth-degree coefficient forces different observational DAG fibers.  Consider an arbitrary strict representation \((G',\nu,\gamma,P')\in\mathcal M_d^\circ\), and denote its root-cluster polynomials by
\[
K_v=z_v\prod_{w\in\operatorname{Ch}_{G'}(v)}
       (1-\gamma_{vw}+\gamma_{vw}K_w).
\]
Reverse-topological induction shows that every \(K_v\) has nonnegative coefficients, has zero constant coefficient, and is divisible by \(z_v\).  Moreover,
\[
[z_v]K_v=\prod_{w\in\operatorname{Ch}_{G'}(v)}(1-\gamma_{vw})>0,
\tag{3}
\]
where positivity uses \(0<\gamma_{vw}<1\) and finiteness of the DAG.  If \(i\to j\) is an edge, selecting the term \(\gamma_{ij}K_j\), its positive singleton term from (3), and the constant term \(1-\gamma_{ih}\) from every other child factor gives
\[
[z_iz_j]K_i\ \geq\
\gamma_{ij}[z_j]K_j
\prod_{h\in\operatorname{Ch}_{G'}(i)\setminus\{j\}}(1-\gamma_{ih})>0.
\tag{4}
\]
Conversely, suppose \([z_iz_j]K_a>0\).  Divisibility by \(z_a\) implies \(a\in\{i,j\}\).  If \(a=i\), a nonconstant term must have been selected from at least one child factor of \(i\).  Every such selected \(K_h\) is divisible by its child's distinct label \(z_h\); exact support \(\{i,j\}\) therefore forces \(h=j\), and hence \(i\to j\).  The case \(a=j\) similarly forces \(j\to i\).  Since \(\nu_a>0\) and all root-cluster coefficients are nonnegative, the log-PGF sum has no cancellation.  We have proved the observable adjacency criterion
\[
[z_iz_j]\bigl(\Psi_{P'}-\Psi_{P'}(0)\bigr)>0
\quad\Longleftrightarrow\quad
\{i,j\}\text{ is an edge of the skeleton of }G'.
\tag{5}
\]

Apply (5) to a graph \(G'\in\mathcal D(P^{(r)})\), using \texttt{def:observational-fiber}.  Equation (1) has a positive coefficient for every pair in \(L\), so on these three vertices \(G'\) has the complete skeleton.  Every acyclic orientation of that skeleton has a unique topological order, say
\[
a\longrightarrow b\longrightarrow c,
\qquad a\longrightarrow c,
\]
and \(c\) is its unique sink.  Write the three strict edge probabilities as \(p=\gamma_{ab}\), \(q=\gamma_{bc}\), and \(t=\gamma_{ac}\).  The same root-cluster recursion gives
\[
K_c=z_c,
\qquad
K_b=z_b(1-q+qz_c),
\qquad
K_a=z_a(1-p+pK_b)(1-t+tz_c).
\]
Only \(K_a\) can contain all three labels, and direct coefficient extraction yields
\[
\left[z_az_bz_c^2\right]
\bigl(\Psi_{P'}-\Psi_{P'}(0)\bigr)=\nu_a pqt>0.
\tag{6}
\]
The displayed recursion also shows that the coefficients of \(z_a^2z_bz_c\) and \(z_az_b^2z_c\) are zero: \(z_a\) occurs only as the leading factor of \(K_a\), while \(z_b\) occurs in \(K_a\) only through the single factor \(K_b\); acyclicity prevents either label from reappearing downstream.  Thus among the three monomials \((\prod_{v\in L}z_v)z_q\), the positive one uniquely names the sink \(c\).

Equations (2) and (6) force \(c=r\) in every \(G'\in\mathcal D(P^{(r)})\).  There are exactly two acyclic orientations of the complete triangle with sink \(r\), according as \(s\to m\) or \(m\to s\).  Both occur: the first is the displayed construction, and interchanging the roles of \(s\) and \(m\), together with their assigned innovation rates and incident edge probabilities in \texttt{def:triangle-witness}, leaves the symmetric expression (1) unchanged.  Equality of the resulting log-PGFs implies equality of their PGFs; coefficient extraction from the PGFs gives equality of every count-cell probability.  Hence
\[
\mathcal D(P^{(r)})=
\left\{
\begin{array}{l}
s\to m,\ s\to r,\ m\to r,\\
m\to s,\ m\to r,\ s\to r
\end{array}
\right\},
\tag{7}
\]
where each line denotes its edge set.  In particular the three fibers have different, indeed mutually exclusive, triangle sinks.

Finally let \(d>3\), put \(W=V\setminus L\), and append, identically for all three laws, mutually independent coordinates \(X_w\sim\operatorname{Pois}(\lambda_w)\) with \(\lambda_w>0\), isolated from the triangle.  Independence and the Poisson PGF give
\[
\widetilde\Psi^{(r)}(z)
=\Psi^{(r)}(z_L)+\sum_{w\in W}\lambda_w(z_w-1).
\tag{8}
\]
Thus the constants and every coefficient through degree three still agree for the three choices of \(r\), while (2) remains unchanged.  All pair coefficients involving a vertex of \(W\) vanish; criterion (5) therefore forces every such vertex to be isolated in every compatible DAG.  The restriction to \(L\) is consequently exactly (7), so the three extended laws again have different fibers.  This proves the claimed sharpness for every \(d\geq3\).
\end{proof}
\par\smallskip\noindent \textit{Justification.} The explicit strict witness proves that degree four is minimal rather than merely sufficient. \textit{Gap.} The cited PB-SCM papers exhibit direction asymmetry but do not provide two strict laws agreeing through every degree-three log-PGF coordinate. \textit{Consumer.} Any PGF procedure claiming universal recovery can be tested against the precise fourth-degree obstruction.

\begin{theorem}[thm:coefficient-confidence]\label{thm:coefficient-confidence}
For every distribution \(P\) on \(\mathbb N_0^V\) with \(p_0>0\), \(T_k((p_x)_{x\in I_4})=c_k\) for \(1\le|k|\le4\), and \(\Pr_P\{t_4(P)\in C_n\}\ge1-\eta\). The inequality is finite-sample and uniform over \(\mathcal M_d^\circ\), without a lower bound on \(p_0\).
\end{theorem}
\begin{proof}
Let
\[
 G_P(z):=\mathbb E_P\!\left[\prod_{i\in V}z_i^{X_i}\right]
 =\sum_{x\in\mathbb N_0^V}p_xz^x.
\]
Because \(p_0>0\), in the ring of formal power series we may write
\[
 G_P(z)=p_0(1+R(z)),\qquad
 R(z)=\sum_{x\ne0}\frac{p_x}{p_0}z^x.
\]
Every monomial of \(R\) has positive total degree.  Therefore, for a multi-index \(k\) with \(1\le |k|\le4\), only the powers \(R^\ell\) with \(1\le\ell\le |k|\) can contribute to \([z^k]\log(1+R)\), and every contributing \(x_a\) satisfies \(|x_a|\le4\).  The formal identity \(\log(1+R)=\sum_{\ell\ge1}(-1)^{\ell+1}R^\ell/\ell\) consequently gives
\[
\begin{aligned}
 c_k
 &=[z^k]\bigl(\log G_P(z)-\log p_0\bigr)\\
 &=\sum_{\ell=1}^{|k|}\frac{(-1)^{\ell+1}}{\ell}
   \sum_{\substack{x_1+\cdots+x_\ell=k\\x_a\ne0}}
   \prod_{a=1}^{\ell}\frac{p_{x_a}}{p_0}
 =T_k((p_x)_{x\in I_4}).
\end{aligned}
\tag{11}
\]
The inner sum is over ordered tuples, exactly as in the definition of \(T\).  Thus no analytic radius-of-convergence or omitted higher-count cell is involved.

It remains to prove simultaneous coverage.  Put \(L=|I_4|\), and for fixed \(x\in I_4\) set \(Y_b=\mathbf 1\{X^{(b)}=x\}\).  Under independent sampling, the \(Y_b\) are independent Bernoulli variables with mean \(p_x\).  For completeness, define
\[
 f(\lambda)=\log\mathbb E\exp\{\lambda(Y_b-p_x)\}.
\]
Under exponential tilting, \(f''(\lambda)\) is the variance of a Bernoulli variable, and hence \(f''(\lambda)\le1/4\).  Since \(f(0)=f'(0)=0\), integration twice gives
\[
 f(\lambda)\le\frac{\lambda^2}{8}\qquad(\lambda\in\mathbb R).
\]
Independence and Markov's inequality then imply, for \(r>0\),
\[
\begin{aligned}
 \Pr\{\widehat p_x-p_x\ge r\}
 &\le\inf_{\lambda>0}
   \exp\left(-\lambda nr+\frac{n\lambda^2}{8}\right)
 =e^{-2nr^2},\\
 \Pr\{p_x-\widehat p_x\ge r\}&\le e^{-2nr^2},
\end{aligned}
\]
where the first infimum is attained at \(\lambda=4r\), and the second line applies the same argument to \(-Y_b\).  A union bound at \(r=\epsilon_n\) yields
\[
 \Pr_P\!\left\{
   \max_{x\in I_4}|\widehat p_x-p_x|>\epsilon_n
 \right\}
 \le 2L e^{-2n\epsilon_n^2}=\eta.
\tag{12}
\]
On the complementary event, the vector \(q=(p_x)_{x\in I_4}\) is feasible for \(Q_n\): it is coordinatewise nonnegative, its truncated sum is at most one, \(q_0=p_0>0\), and it satisfies every interval constraint.  By (11), \(T(q)=t_4(P)\), whence \(t_4(P)\in C_n\).  Combining this with (12) proves
\[
 \Pr_P\{t_4(P)\in C_n\}\ge1-\eta.
\]
The bound (12) is distribution-free and contains no division by a lower bound for \(p_0\).  It is therefore uniform over \(\mathcal M_d^\circ\), although the random set \(C_n\) may be unbounded when its feasible values of \(q_0\) approach zero.
\end{proof}
\par\smallskip\noindent \textit{Justification.} The finite rational log-series identity converts bounded cell concentration into honest simultaneous coefficient coverage. \textit{Gap.} Xiang et al. use bootstrap rank tests and Nakamura and Perez-Abreu study empirical-PGF inference, but neither constructs this denominator-safe finite coefficient region. \textit{Consumer.} The coefficient region is the input for graph-set inference in sparse football and shopping count tables.

\begin{theorem}[thm:outer-fiber-confidence]\label{thm:outer-fiber-confidence}
The outward inverse satisfies \(\Pr_P\{\mathcal D(P)\subseteq\mathcal D_n\}\ge1-\eta\) uniformly over \(\mathcal M_d^\circ\). If the reporting rule outputs no arrow when \(\mathcal D_n=\varnothing\) and otherwise reports only arrows present with the same orientation in every \(G\in\mathcal D_n\), then the unconditional probability that it reports any arrow not compelled throughout \(\mathcal D(P)\) is at most \(\eta\).
\end{theorem}
\begin{proof}
Let \(\mathcal E_n=\{t_4(P)\in C_n\}\).  The already established thm:coefficient-confidence, under ass:iid-sampling and \(p_0>0\), gives
\[
 \Pr_P(\mathcal E_n)\ge1-\eta
\tag{15}
\]
uniformly over \(\mathcal M_d^\circ\).  (For a strict PB-SCM, \(p_0>0\): all finitely many independent Poisson innovations vanish with probability \(\exp(-\sum_i\mu_i)>0\), and then the structural recursion gives \(X=0\).)

On \(\mathcal E_n\), take any \(G\in\mathcal D(P)\).  By def:observational-fiber there is a strict parameter point on \(G\) inducing \(P\), so \(G\in\mathcal D(t_4(P))\).  Equivalently, thm:degree-four-completeness identifies \(\mathcal D(t_4(P))\) with the complete law fiber \(\mathcal D(P)\).  Since \(t_4(P)\in C_n\), def:outer-fiber gives
\[
 G\in\mathcal D(t_4(P))
 \subseteq\bigcup_{t\in C_n}\mathcal D(t)=\mathcal D_n.
\]
This holds for every \(G\in\mathcal D(P)\); together with (15), it proves
\[
 \Pr_P\{\mathcal D(P)\subseteq\mathcal D_n\}\ge1-\eta
\tag{16}
\]
uniformly over the strict class.

The reporting rule emits nothing when \(\mathcal D_n=\varnothing\).  Otherwise, suppose (16)'s inclusion event holds and the rule emits \(i\to j\).  By definition of the rule, \(i\to j\) belongs to every member of \(\mathcal D_n\), hence, by the inclusion, to every member of \(\mathcal D(P)\); it is therefore compelled throughout the population fiber.  Consequently
\[
 \{\text{at least one arrow not compelled in }\mathcal D(P)\text{ is reported}\}
 \subseteq
 \{\mathcal D(P)\nsubseteq\mathcal D_n\}.
\]
Taking probabilities and using (16) bounds the left side by \(\eta\), uniformly over \(\mathcal M_d^\circ\).  This is unconditional familywise control; no conditioning on nonemptiness and no unsupported division by its probability is used.
\end{proof}
\par\smallskip\noindent \textit{Justification.} Inference targets the full non-singleton population fiber and therefore avoids false certainty from graph-then-report workflows. \textit{Gap.} Strieder et al. invert tests for linear-model effects and Wang et al. give asymptotic ordering sets; neither covers every graph in a strict PB-SCM population fiber with finite-sample probability. \textit{Consumer.} The football-event and shopping-mall analyses can report only causal arrows supported simultaneously across structural and sampling uncertainty.

\begin{openendedquestion}[oeq:polynomial-delay-inversion]\label{oeq:polynomial-delay-inversion}
Does there exist an algorithm which, for every fixed rational confidence instance \((n,\eta,\widehat p)\), returns an output list \(O\) and certificate family \(\kappa\) satisfying \(\mathsf A_{\mathrm{out}}(n,\eta,\widehat p;O,\kappa)\), terminates with delay polynomial in the input bit length and \(d\) between successive output source-forest classes, and thereby preserves the coverage conclusion of \(\text{thm:outer-fiber-confidence}\)?
\end{openendedquestion}
\par\smallskip\noindent \textit{Justification.} The population geometry suggests an exact output-sensitive alternative to generic semialgebraic elimination, but existence of an algorithm satisfying the certified handle specification with polynomial delay is genuinely open. \textit{Gap.} Existing causal-model inversion papers do not exploit a source-forest quotient, and the PB-SCM papers provide point algorithms rather than certified outward enumeration. \textit{Consumer.} A positive answer would make honest full-fiber confidence reporting computationally usable beyond very small event systems.

\section*{Technical internal limitation (diagnostic only)}

The strict interior is essential. At zero thinning an edge is observationally absent, at zero innovation a source can disappear, and at deterministic thinning the positive failure patterns and decoder denominators can fail. The coefficient confidence region may be unbounded when the data do not bound \(p_0\) away from zero. These facts preclude automatic uniform singleton recovery on the unrestricted class and are not repaired by imposing an undeclared compact parameter box.

\section{Honest scope}

The formal claims concern finite labeled causally sufficient DAGs, independent positive Poisson innovations, and edge probabilities strictly between zero and one. As a required sanity reduction, when \(d=2\) and one edge is present, the source-forest theorem yields exactly the two orientations and the transformation in Qiao et al.; when the skeleton is empty, it yields the unique empty DAG. No automatic faithfulness theorem, latent-confounded extension, boundary-fiber theorem, numerical stability bound, or unrestricted uniform singleton-consistency result is claimed. The arbitrary-DAG repeated-label genealogy lemma is the central proof obligation. Exact polynomial-delay outward inversion is posed as an open-ended question rather than asserted.

\begin{thebibliography}{99}

  \bibitem{qiao2024causal} Jie Qiao, Yu Xiang, Zhengming Chen, Ruichu Cai, and Zhifeng Hao (2024). Causal Discovery from Poisson Branching Structural Causal Model Using High-Order Cumulant with Path Analysis. Proceedings of the AAAI Conference on Artificial Intelligence 38, 20524--20531.

  \bibitem{xiang2024identifiability} Yu Xiang, Jie Qiao, Zhefeng Liang, Zihuai Zeng, Ruichu Cai, and Zhifeng Hao (2024). On the Identifiability of Poisson Branching Structural Causal Model Using Probability Generating Function. Advances in Neural Information Processing Systems 37, 11664--11699.

  \bibitem{ghassami2020characterizing} AmirEmad Ghassami, Alan Yang, Negar Kiyavash, and Kun Zhang (2020). Characterizing Distribution Equivalence and Structure Learning for Cyclic and Acyclic Directed Graphs. Proceedings of Machine Learning Research 119, 3494--3504.

  \bibitem{gao2026causal} Penggang Gao, Ming Cai, and Hisayuki Hara (2026). Causal DAG Identification for Count Data via Poisson Thinning Structural Equation Models. arXiv:2609.06098.

  \bibitem{qiao2026latent} Jie Qiao, Zihuai Zeng, Ruichu Cai, Zhengming Chen, and Zhifeng Hao (2026). On the Identifiability of Poisson Branching Structural Causal Model Under Latent Confounding. Proceedings of the Forty-Third International Conference on Machine Learning.

  \bibitem{strieder2021confidence} David Strieder, Tobias Freidling, Stefan Haffner, and Mathias Drton (2021). Confidence in Causal Discovery with Linear Causal Models. Proceedings of Machine Learning Research 161, 1217--1226.

  \bibitem{wang2026confidence} Y. Samuel Wang, Mladen Kolar, and Mathias Drton (2026). Confidence Sets for Causal Orderings. Journal of the American Statistical Association, published online in 2025.

  \bibitem{nakamuraPerezAbreu1993poisson} Miguel Nakamura and Victor Perez-Abreu (1993). Use of an Empirical Probability Generating Function for Testing a Poisson Model. Canadian Journal of Statistics 21, 149--156.

  \bibitem{nakamuraPerezAbreu1993overview} Miguel Nakamura and Victor Perez-Abreu (1993). Empirical Probability Generating Function: An Overview. Insurance: Mathematics and Economics 12, 287--295.

  \bibitem{steutelVanHarn1979} Fred W. Steutel and Klaas van Harn (1979). Discrete Analogues of Self-Decomposability and Stability. Annals of Probability 7, 893--899.

  \bibitem{alOshAlzaid1987} M. A. Al-Osh and A. A. Alzaid (1987). First-Order Integer-Valued Autoregressive Process. Journal of Time Series Analysis 8, 261--275.

\end{thebibliography}

\end{document}